\section{Limitations and Open Questions}
\label{sec:limitations}

\textbf{Stated plainly, because each one shapes what we would do next.}

\emph{No forgetting audit yet.} We have not measured general-capability
regression (e.g.\ an MMLU subset) after DAPT; the 8\% replay mix follows
ChipNeMo practice but is unverified here. This is the first gap we would
close.

\emph{Single seed, single platform.} Every number is one training run on one
platform; the ablation ladder's steps are large (+37.6 points for the recipe
change) but error bars would still become necessary before any causal claim
about smaller effects.

\emph{Bundled recipe components.} Full coverage, template count, reversal, and
whole-table narratives shipped together in the final run; their individual
contributions are not isolated. The interference explanation
(Finding~5) in particular rests on one observation and needs a controlled
similarity$\times$exposure$\times$context experiment.

\emph{L3 does not differentiate.} Register naming is descriptive by design, so
description$\leftrightarrow$name multiple choice is partially solvable by
semantic matching; all models score $\ge$89\% there even after leak
filtering.

\emph{Remaining model weaknesses.} Numeric range membership (which region
contains address $X$: 11/27) resists declarative augmentation; it needs
reasoning over the map, suggesting either chain-style augmentation or an SFT
stage; issue-history recall remains weak for all models.

\emph{No RAG comparison.} This study measures what training injects into
weights; a deployment-oriented comparison against retrieval (and against
retrieval \emph{plus} our adapted model) is the obvious next experiment, and
we expect the hybrid to win.
